\section{The finite marked-pair identity}
\label{sec:06-magnetic-algebra}

The purpose of this section is to express the positive two-mark
observable as an exactly normalized finite contour sum. The spin
contraction selects one polygon, the dual covectors resolve its
ordinary sites, and residue cancellation produces the contour identity.

We retain the periodic vacuum \(U_N\), its polynomial normalization \(h_N\),
and the constants
\[
 \lambda=\frac{\pi}{8},\qquad \beta=2\lambda=\frac{\pi}{4},\qquad
 q=e^{i\beta},\qquad k=q^3,\qquad
 c=\cos(3\lambda),\qquad d=\cos(2\lambda),\qquad
 \gamma=\frac{\pi}{3\beta}.
\]
The cyclic order of the sites in this section is
\([p,X,p',Y]\).  The two marked objects \(p,p'\) are sites, or ports,
and \(X,Y\) are the ordinary sites.  Set
\[
 s=|X|,\qquad b=|Y|,\qquad n=s+b,\qquad N=n+2,\qquad
 \sigma_i=\begin{cases}1&i\in X,\\-1&i\in Y.\end{cases}
\]
Throughout, \(n\) is even.  At physical parameters let
\(\mathcal C(p,p')\) be the sum of the weights of individual cylinder
polygons visiting both marked ports.  Each polygon has its original
vertex-activity weight.  Below, the symbol \(\rho\) will instead denote
the probability density used to center the contour gas.

\subsection{A bilinear spin contraction}

Use the three-dimensional spin space with basis indexed by \(0,+1,-1\),
where \(0\) denotes a vacancy. An ordered cup or cap has coefficient
\(k^h\) on the spins \(h,-h\) for \(h\in\{+1,-1\}\),
and a straight strand preserves its spin.  These assignments represent
composition of diagrams by ordinary matrix multiplication: two
successive bends contribute \(k^h k^{-h}=1\), and a closed unmarked loop
contributes \(k^2+k^{-2}=0\).  All transposes and contractions here are
bilinear, without complex conjugation.  Each local box preserves total
spin.

For the star list write \(u_i^*=3\lambda-u_i\) and
\(U^*=U_N((u_i^*)_i)\).  The star on this symbol thus denotes a
parameter substitution, not an adjoint.

\begin{lemma}[Marked spin contraction]
\label{lem:06-spin-contraction}
In the linear order starting at \(p\),
\begin{equation}
 \mathcal C(p,p')
 =
 \sum_{\mathbf h\in\{0,\pm1\}^{X\cup Y}}
 k^{-2\sum_{i\in X}h_i}\,
 U_{+,\mathbf h_X,-,\mathbf h_Y}\,
 U^*_{-,\mathbf h_X,+,\mathbf h_Y}.
 \label{eq:06-spin-contraction}
\end{equation}
\end{lemma}

\begin{proof}
First fix a pair of disk matching diagrams in the two half-cylinder
vacua.  Draw their common cut vertically upwards, with the sites in the
stated order, and their matching arcs on the two sides of this line.
The ket occupies the left side.  Orient a left arc from \(+\) to \(-\),
and a right arc from \(-\) to \(+\).
Away from the marked ports the spin contractions concatenate directed
arcs.  At each marked port the two prescribed spins instead give the
two endpoints of directed paths.  Thus a nonzero assignment forces two
disjoint simple directed arcs from \(p\) to \(p'\), constituting a
single marked loop.  Their directions are fixed, with no remaining
factor of two.

Every other loop avoids the marked loop and cannot separate its two
marked points.  Its algebraic crossing number through the cut interval
\(X\) is therefore zero.  The twist in
\eqref{eq:06-spin-contraction} does not change its orientation sum,
which is the ordinary loop value \(k^2+k^{-2}=0\).
It remains to compute the marked loop's phase.

Let \(W\) be the sum of the tangent turns along its two directed arcs,
and let \(F_X\) be their total number of westward minus eastward
crossings of \(X\).  The cup and cap coefficients contribute
\(k^{-W/\pi}\), and the inserted twist contributes \(k^{-2F_X}\).
We claim that \(W=-2\pi F_X\).
One may smooth the semicircular matching arcs at their joins.
For a simple regular arc \(\Gamma\) from \(p\) to \(p'\), consider the
direction of the chord \(\Gamma(t)-\Gamma(s)\) on the triangle \(s\le t\).
It extends continuously to its diagonal by the tangent direction.
Since the triangle is simply connected, its argument has a continuous
lift.  Comparing its three boundary sides shows that the tangent turn
equals the sum of the argument changes of
\(\Gamma-p\) and \(p'-\Gamma\), with tangent limits at the vanishing
endpoints.

For \(\Gamma-p\), put the argument cut downwards.  There are no sites
below \(p\), so the matching drawing has no crossings of this cut.
The initial angles on the two arcs are \(0,\pi\), and the final angles
are both \(\pi/2\); their summed change is zero.
For \(p'-\Gamma\), use the upward cut and the range
\((-3\pi/2,\pi/2)\).  The sums of the initial and final one-sided
angle limits are both \(-\pi\).  Each algebraic crossing of the cut
contributes \(-2\pi\), giving total change \(-2\pi F_X\).
This proves the claim, so the marked loop has phase one.

The disk matching drawing is also valid when the resulting polygon
winds around the cylinder: no vacuum strand runs to an infinite end.
Finally the absolute convergence of the physical half-cylinder vacuum
sums allows the diagram-by-diagram calculation to be summed.
\end{proof}

\subsection{The finite algebra of the marked pair}

Put \(a_i=2iu_i\), \(x_i=e^{a_i}\), and
\[
 P_j(v)=-2\sinh\!\left(\frac{v+ij\beta}{2}\right).
\]
We retain the centered normalization \(\widetilde h_j\) from
Section~\ref{sec:reference}. The finite gas integral \(I\) has the
following factors and contours. Within a group, its site--site,
particle--site and particle--particle factors are respectively
\(B_s,F_s,E_s\); across groups they are \(B_c,F_c,E_c\):
\[
\begin{array}{lll}
 B_s=P_2P_{-2}P_3P_{-3},&
 F_s=(P_1P_{-1})^{-1},&
 E_s=P_0^2P_2P_{-2}F_s,\\
 B_c=P_2^2P_3^2,&F_c=P_1^{-2},&
 E_c=P_0^2P_2^2F_c .
\end{array}
\]
The within-group factors are even. Across groups, evaluate each
factor at the \(X\) coordinate minus the \(Y\) coordinate.
For each group there are main singletons, far singletons, and far
doubles, of elementary multiplicities $1,1,2$. A double has one real
integration variable and two fixed imaginary shifts. After a common
shift, the ordinary \(X\) sites and both marks have heights \(0\) or
\(\beta\), and the ordinary \(Y\) sites have height zero. Real site
offsets are first generic and distinct. The horizontal contours have
the following heights, in units of \(\beta\):
\[
\begin{array}{c|cc}
 &X&Y\\ \hline
 \text{main singleton}&1/2&0\\
 \text{far singleton}&25/8&-7/2\\
 \text{far double}&(3+\varepsilon,4+\varepsilon)&(-3,-4).
\end{array}
\]
Here \(\varepsilon>0\) is fixed and sufficiently small, and all
horizontal lines are oriented from left to right.
Lemma~\ref{lem:06-contours} will justify their deformation and
confluence, including the curved main \(X\) contour used later.
Multiply all
site--site factors, all particle--site factors and all interactions
between distinct elementary particles, omitting the internal pair of
each double. Integrate each real variable with measure $dv/(2\pi)$
(the curved main differential is used after deformation). For species
counts $n_\sigma,f_\sigma,d_\sigma$, include the activity and factorial
factor
\[
 \prod_{\sigma=\pm1}
 \frac{p_\sigma^{n_\sigma}(-p_\sigma)^{f_\sigma}
             (4c^2p_\sigma^2)^{d_\sigma}}
      {n_\sigma!f_\sigma!d_\sigma!}.
\]
Sum only through total elementary degree at most $n$. This defines
$I$ as a finite polynomial of contour integrals, not as an infinite
activity expansion.

Only \(X\cup Y\) supplies sites to this gas.  A double is always
interpreted as its two elementary particles when a product over
particles is written.  The marked-site insertion is
\begin{equation}
 \mathcal T=
 \prod_{a\in\{a_p,a_{p'}\}}
 \left[
 \prod_{i\in X\cup Y}
 P_{2\sigma_i}(a_i-a)P_{3\sigma_i}(a_i-a)
 \prod_{v\ {\rm elementary}}
 \frac{P_{-4\sigma_v}(v-a)}{P_{\sigma_v}(v-a)}
 \right].
 \label{eq:06-insertion}
\end{equation}
The two marks are treated separately, even when their spectral
parameters agree.  Write \(I(\mathcal T)\) for the integral with this
additional factor. These finite integrals are absolutely convergent.
Indeed, if a cluster of elementary degree $r\ge1$ in a term of total
degree $m\le n$ is translated to either infinity, the site factors
supply $e^{-nr|v|}$ and its interactions with the other particles
supply $e^{r(m-r)|v|}$. Thus its tail is
$O(e^{-r(n-m+r)|v|})$. The marked insertion ratios are bounded at
both ends. Ordering the real integration coordinates and applying
this estimate to each separating gap proves joint integrability;
after combining the factors of each double, the specified contours
avoid every pole.
Only degrees $n$ and $n-1$ will be used, respectively, for marked
polygons and one arc.

\begin{theorem}[Marked-pair contour identity]
\label{thm:06-marked-identity}
At the physical parameters, and by continuation at regular complex
parameters on the contours of Lemma~\ref{lem:06-contours},
\begin{equation}
 \widetilde h_N^2\mathcal C(p,p')
 =(2c)^2(2d)^n
 \sum_{M=-\min(s,b)}^{\min(s,b)}
 (-1)^{sb+M}
 [p_+^{\,s-M}p_-^{\,b+M}]I(\mathcal T).
 \label{eq:06-marked-identity}
\end{equation}
Here \(p_+,p_-\) are formal activities, unrelated to the marked port
\(p\).
\end{theorem}

We prove the finite algebra first.  Its two ingredients are stated
separately to make clear where the normalization in
\eqref{eq:06-marked-identity} enters.

\begin{lemma}[A complete pair of dual covector families]
\label{lem:06-dual-covectors}
Fix one ordinary group of sign \(\sigma\), containing \(g\) sites.
For each \(\ell\in\{1,0,-1\}^g\), construct a covector \(D_\ell\) as
follows.  Split a site with \(\ell_j=0\) by \(S_2\) into arguments
\(u_j-\lambda,u_j+\lambda\), with relative signs \(+,-\).
An unsplit site retains its sign \(\ell_j\).
Stably braid the resulting roots into pluses followed by minuses,
and project their spins to \(\sigma\) times their relative signs.
On the star list, use signs \(-,+\) in each split, stably sort minuses
before pluses, and project in the same relative-spin convention;
call the resulting covector \(D_\ell^*\).

At generic arguments,
\begin{equation}
 D_\ell(D_r^*)^{\mathsf t}
 =\mathbf 1_{\ell=r}\,d_\ell,\qquad
 d_\ell=\prod_{\substack{z\ {\rm plus\ root}\\w\ {\rm minus\ root}}}
 F(u_w-u_z).
 \label{eq:06-biorthogonality}
\end{equation}
Here \(F\) is the local cap--cup box coefficient, rather than a
site--particle gas factor.  In particular \(d_\ell\ne0\) generically,
and these \(3^g\) covectors give a complete bilinear resolution of the
identity.  Each covector has input charge
\(\sigma\sum_j\ell_j\).
\end{lemma}

\begin{proof}
The root lists associated with a site and its star correspond under
\(u\mapsto3\lambda-u\) with the same relative signs.  The splitting
relation transports an exchange of two underlying sites to the
exchange of their two blocks of roots.  After sorting, this agrees
with the original sorting up to exchanges within a fixed sign class.
The latter act on the projected spin by the nonzero scalar \(E\).
The Yang--Baxter, unitarity, and disjoint-commutation identities show
that different reduced sequences for the same permutation agree.
This comparison may first be made with generic roots: none of the
braids being compared exchanges the two roots of the same split,
so specialization to a split is legitimate.
The transports of the real and star lists cancel in their bilinear
pairing by symmetry and unitarity.  Consequently the vanishing or
nonvanishing of \(D_\ell(D_r^*)^{\mathsf t}\) is unchanged when the same
permutation is applied to the underlying sites of both lists.

If \(\ell_j>r_j\), put site \(j\) first.  A leading plus root on one
side and a leading minus root on the other are then untouched by
sorting.  At least one of these two sites is unsplit.  It forces its
input spin, whereas a leading root of a split can arise only from
vacancy or from the input spin of that root.  These conditions are
incompatible.  If \(\ell_j<r_j\), use the analogous argument with
site \(j\) last.  This proves the off-diagonal vanishing.

On the diagonal, replace each \(S_2S_2^{\mathsf t}\) by
\(R(2\lambda)\).  The remaining network sorts the minuses-first list
into the pluses-first list.  Each opposite-sign pair crosses once,
and no equal-sign pair needs to cross.  Follow the first plus through
all the minuses.  At any crossing it can only transfer charge to the
new minus; because its final spin is again the maximal plus spin,
it cannot make such a transfer.  Each crossing therefore has the
interchanged-spin coefficient \(F\).  Removing this strand and
repeating proves the product in \eqref{eq:06-biorthogonality}.
It is nonzero at generic arguments.  Since there are exactly \(3^g\)
covectors in a \(3^g\)-dimensional space, the two families are bases.
Explicitly their identity resolution is
\[
 \operatorname{Id}
 =\sum_{\ell}\frac{(D_\ell^*)^{\mathsf t}D_\ell}{d_\ell}.
\]
\end{proof}

\begin{lemma}[Two occupied spin blocks]
\label{lem:06-two-blocks}
For a list of \(2m\) sites in linear order and \(d_0\in\{1,-1\}\),
\begin{equation}
 \widetilde h_{2m}\,U_{d_0^m(-d_0)^m}
 =(2c\,k^{d_0})^m
 \prod_{\substack{i<j\\i,j\ {\rm in\ the\ same\ block}}}
 P_2(a_i-a_j)P_3(a_i-a_j).
 \label{eq:06-two-blocks}
\end{equation}
Also \(\widetilde h_j\) is unchanged on replacing its entire list by
the star list.
\end{lemma}

\begin{proof}
Within either equal-spin block, fusion at the ascending differences
\(2\lambda,3\lambda\) makes the amplitude vanish.  For nonadjacent
sites move them together by braid covariance; equal-spin exchanges
are scalars, so the same zeros hold.  For any occupied slot the
centered polynomial degree bound has exactly the degree supplied by
these two factors for each of the other \(m-1\) sites in its block.
Thus their product exhausts the degree in every slot, leaving a
constant depending only on \(m,d_0\).

Determine the constant by deleting the adjacent pair across the two
blocks at multiplicative arguments \(x,y=q^3x\).
In the centered deletion quotient
\eqref{eq:vacuum-centered-deletion}, each spectator factor is the
negative of its removed same-block
\(P_2P_3\) factor.  There are \(2m-2\) spectators, so these signs
multiply to one.  The inserted cup contributes \(k^{d_0}\).
The constant therefore satisfies
\(C_m=(2c\,k^{d_0})C_{m-1}\), with \(C_0=1\), proving
\eqref{eq:06-two-blocks}.  Finally reciprocity of \(h_j\), together
with its homogeneity, gives the asserted equality of the centered
normalizations.
\end{proof}

\begin{proof}[Proof of Theorem \ref{thm:06-marked-identity}: finite part]
Resolve the contractions on \(X\) and \(Y\) with
Lemma \ref{lem:06-dual-covectors}.  The twist acts by a scalar on each
charge subspace.  Total spin conservation, including the marked
spins, leaves exactly
\[
 \sum_{i\in X}\ell_i=\sum_{i\in Y}\ell_i.
\]
Splitting and braid covariance of the normalized vacuum identify the
resulting projected amplitudes with amplitudes on the sorted root
lists.  Write \(X_\pm,Y_\pm\) for the roots with the indicated relative
signs.  Move the ket tail \(Y_-\) to the front, and the bra tail
\(Y_+\) to the front.  The ket blocks are
\([Y_-,p,X_+]\), \([X_-,p',Y_+]\), while the bra blocks are
\([Y_+,p,X_-]\), \([X_+,p',Y_-]\).
Both pairs have equal block sizes.

Cyclic covariance is an identity of matching diagrams.  In the spin
embedding, moving a tail to the front reverses the endpoint order
exactly for cups joining the tail to its complement. If the tail
endpoint has spin \(h\), its old coefficient divided by its new one
is \(k^{-2h}\). Cups internal to the tail have total spin zero.
Thus a moved tail of total spin \(H\) changes the amplitude, old
order relative to new order, by \(k^{-2H}\).
For the two tails above the product is
\(k^{2\sum_Y\ell_i}\), canceling the original twist.
The phases in the two constants of
\eqref{eq:06-two-blocks} also cancel.  We next record all remaining
factors, including the divisors from the change of normalization.

At an ordinary site of label \(\ell\), encode each root by its
relative sign \(e\) and its shift \(h\), so that its log coordinate is
\(a_i+ih\beta\):
\[
 \mathcal R(1)=\{(1,0)\},\qquad
 \mathcal R(-1)=\{(-1,0)\},\qquad
 \mathcal R(0)=\{(1,-1),(-1,1)\}.
\]
Clearing \(\widetilde h_N^2\) on the original list, instead of the
expanded list, divides by the squares of
\eqref{eq:vacuum-centered-merge}. Since \((2d)^2=2\) and
\(P_4(\log(z/x))P_{-4}(\log(z/x))=(z+x)^2/(zx)\), each
expanded site contributes a scalar divisor \(2\), and a spectator
divisor \(P_4P_{-4}\).  For two ordinary sites this divisor occurs
once if exactly one is expanded; if both are expanded, it occurs
at shifts \(-1,0,1\) of their log difference.
The diagonal \(d_\ell\) also contains, within every expansion,
\(F(2\lambda)=(2c)^2/2\).  If \(z\) sites are expanded, the total
number of occupied roots including the marks is \(N+z\).
The two block constants contribute \((2c)^{N+z}\), and the scalar
divisors contribute \((2c)^{2z}\).  Their quotient is therefore
\[
 (2c)^{N-z}=(2c)^2(2c)^{\sum_i\ell_i^2}.
\]

For an ordered ordinary pair with groups \(\sigma,\omega\), multiply
the following factor over
\((e,h)\in\mathcal R(\ell)\), \((f,g)\in\mathcal R(r)\), evaluating it
at shift \(h-g\) of the log difference:
\[
\begin{array}{c|l}
 \text{condition}&\text{factor}\\ \hline
 \sigma=\omega,\ f=e & P_2P_3P_{-2}P_{-3}\\
 \sigma=\omega,\ f=-e & P_{-2e}P_{-3e}/(P_eP_0)\\
 \sigma=1,\ \omega=-1,\ f=-e & P_{-2e}^{\,2}P_{-3e}^{\,2}\\
 \text{otherwise}&1 .
\end{array}
\]
A string of length $a$ at a group-$\sigma$ site $a_i$ consists of
particles at $a_i+i\sigma j\beta$, $1\le j\le a$; length zero
means the empty string. Divide by the merge divisors just specified.
The result is precisely
the ordinary mutual string product with lengths
\(1+\sigma\ell,1+\omega r\), multiplied by \(1\) within a group and
\((-1)^{1-\ell r}\) between groups.
Here is an exponent check that also fixes the signs.
Encode \(\prod_j P_j(v)^{m_j}\) by the Laurent polynomial
\(\sum_j m_j\xi^j\); shifting \(v\) by \(ih\beta\) multiplies
this polynomial by \(\xi^h\).
For string lengths \(a,b'\), put
\[
 A=\sum_{h=1}^{a}\xi^{\sigma h},\qquad
 D=\sum_{h=1}^{b'}\xi^{-\omega h}.
\]
The string polynomial is \(B-F(A+D)+EAD\), where
\[
\begin{array}{c|ccc}
 &B&F&E\\ \hline
 \sigma=\omega&
 \xi^2+\xi^{-2}+\xi^3+\xi^{-3}&\xi+\xi^{-1}&
 2+\xi^2+\xi^{-2}-\xi-\xi^{-1}\\
 \sigma=1,\omega=-1&
 2\xi^2+2\xi^3&2\xi&2+2\xi^2-2\xi .
\end{array}
\]
Subtract this polynomial from the root-product polynomial and
subtract \(\xi^4+\xi^{-4}\) at every divisor shift.
The result is zero within either group.  Across groups it is the
following matrix, with rows and columns in the order \(+,0,-\):
\[
 \begin{pmatrix}
 0&b_4+2b_5&2b_5+2b_6\\
 -b_4&b_4+b_3-b_5&b_4+2b_5\\
 0&-b_4&0
 \end{pmatrix},
 \qquad b_j=\xi^{j-8}-\xi^j.
\]
Since \(P_{j-8}=-P_j\), each \(b_j\) contributes a minus sign.
The displayed matrix gives exactly \((-1)^{1-\ell r}\).

A root \((e,h)\) interacts with either marked spectator \(a\)
by \(P_{-2e}P_{-3e}\), at shift \(h\).  Dividing by the merge divisor
when necessary gives, for \(\ell=1,0,-1\), respectively,
\[
 P_{-2}P_{-3},\qquad P_{-3}P_3,\qquad P_2P_3
\]
at \(a_i-a\).  Substituting the strings of lengths \(1+\sigma\ell\)
in \eqref{eq:06-insertion} gives these same three factors.
There is no marked--marked pair factor: the two marks lie in
different blocks in both amplitudes.

Set \(M=-\sum_X\ell_i=-\sum_Y\ell_i\).
The particle totals are \(s-M,b+M\), and the cross signs multiply to
\[
 (-1)^{sb-(\sum_X\ell_i)(\sum_Y\ell_i)}
 =(-1)^{sb-M^2}=(-1)^{sb+M}.
\]
To compute the string weights, move a group-\(\sigma\) main line
toward its far line, in direction \(\sigma\). With measure \(dv/(2\pi)\),
the old integral equals the new one plus \(i\sigma\) times the
sum of enclosed residues. Directly from the \(P_j\) factors,
\begin{equation}
 \operatorname*{Res}_{v=i\sigma\beta}F_s(v)=\frac1{2i\sigma d},
 \qquad
 \operatorname*{Res}_{v=i\sigma\beta}E_s(v)=4i\sigma c^2,
 \qquad F_s(2i\sigma\beta)=-\frac1{2d}.
 \label{eq:06-string-residues}
\end{equation}
Consequently the empty, one-root, and two-root residues are
\[
 r_0=1,\qquad r_1=\frac1{2d},\qquad
 r_2=r_1(-4c^2)\left(-\frac1{2d}\right)=\frac{c^2}{d^2}.
\]
For two identical particles, the two choices of the first root cancel
their factorial \(2!\).
The numbers of strings of lengths zero and two are equal because the
total elementary count is \(n\).  Thus
\((2d)^n\prod_i r_{1+\sigma_i\ell_i}
=(2c)^{\sum_i\ell_i^2}\).
This proves \eqref{eq:06-marked-identity} for the finite string sum.
The following contour argument completes the proof.
\end{proof}

\subsection{The magnetic contours}

It remains to show that the contours defining \(I\) produce exactly
the strings used in the finite calculation, and that they permit
the later deformation of the main \(X\) line.

\begin{lemma}[Admissible contour deformation]
\label{lem:06-contours}
These contours give the string sum in
\eqref{eq:06-marked-identity}.  The main \(X\) contour may further be
replaced by any smooth graph with bounded derivative and heights in
\([\beta/2,\beta)\).  Height \(\beta\) is allowed when all \(X\) sites
have that height.  The identity passes to coinciding real offsets.
\end{lemma}

\begin{proof}
Write \(A=25/8\), \(C'=7/2\).  First move the double heights to
\((A-1,A)\) in \(X\) and \((-C'+1,-C')\) in \(Y\).
Then move the main contours to their far-singleton contours.
The same-group site poles produce strings of lengths \(0,1,2\).
There is room for both roots before the endpoint in each case,
including for a site of height \(\beta\) in \(X\).
An attempted third root is separated by \(2i\sigma\beta\) from
the first; the zero of \(P_{-2\sigma}\) in their interaction cancels
the pole from the second root. Thus no longer string occurs.
The inequality \(A+C'<7\) excludes cross-group singleton poles.
The interaction of a double with another particle has no poles:
up to translation of the argument, the fused factors are
\[
 E_s(v)E_s(v+i\beta)=P_0(v)P_1(v)P_{-2}(v)P_3(v),
 \qquad
 E_c(v)E_c(v+i\beta)=P_0(v)^2P_3(v)^2.
\]
Together with the stated heights, these identities make the double
movements regular. Finally
\(P_{\sigma_v}(v-a)^{-1}\) introduces no crossed pole on any of these
paths.  These observations also exclude poles during the subsequent
deformation of the main \(X\) contour, including interactions of two
moving main particles.

For completeness, the cancellation on the far contours is
coefficientwise in \(p_+,p_-\).  An unpicked main singleton cancels
an original far singleton because their activities sum to zero.
A main particle at \(v\) attaching to an original far singleton at
\(w\) crosses the opposite pole, \(v-w=-i\sigma\beta\).
Its oriented residue is
\(i\sigma\operatorname*{Res}_{z=-i\sigma\beta}E_s(z)=4c^2\),
by \eqref{eq:06-string-residues} and evenness of \(E_s\).
It therefore produces a double canceling the original double:
\[
 4c^2p_\sigma^2+4c^2(-p_\sigma)p_\sigma=0.
\]
The factorial measures give precisely the binomial coefficients
for mixing these two origins of each final species.  Thus only
site-bound strings remain.

There are \(n\) ordinary sites and \(n\) elementary particles in each
coefficient under consideration.  The original confinement argument
therefore gives simultaneous exponential integrability at infinity.
Every factor from the marked insertion is bounded there, and fused
residues have the same confinement.  One may consequently perform
the one-variable deformations and cancellations under all spectator
integrals.  Continuing from generic distinct offsets gives
confluence: the physical transfer and its normalized vacuum are
regular there, while the contour integrals remain absolutely
convergent. Together with the finite calculation above, this completes
the proof of Theorem~\ref{thm:06-marked-identity}.
\end{proof}
